% SPDX-License-Identifier: CC-BY-SA-4.0
% Author: Matthieu Perrin
% Part: <Nom de la partie>
% Section: <Nom de la section>
% Sub-section: <Nom de la sous-section>  % (facultatif, laisser vide si non utilisé)
% Frame: <Titre de la slide>

\begingroup

\begin{frame}{Calculabilité Et Complexité}

  \on[y=14mm, x=-5mm]{
    \begin{tikzpicture}[2Darray, x=40mm, y=10mm]
      \arrayColumn[width=18mm]{\small
        \arrayLine[draw=none, align=right]{\structure{Structure algébrique}}
        \arrayLine[draw=none, align=right]{\structure{Problème fondamental}}
        \arrayLine[draw=none, align=right]{\structure{Applications à la logique}}
        \arrayLine[draw=none, align=right]{\structure{Systèmes de preuves}}
      }
      
      \arrayColumn[header=\structure{Calculabilité}]{\small
        \arrayLine{\textsc{r}, \textsc{re}, \textsc{re-complet}\\  $\leq_m$, $\equiv_m$}
        \arrayLine{\structure{\textsc{Universal}}\\est \textsc{re-complet}}
        \arrayLine{\structure{\textsc{entscheidungsproblem}} \\ est \textsc{re-complet}}
        \arrayLine{Absence de preuves\\ (Gödel)}
      }
      
      \arrayColumn[header=\structure{Complexité}]{\small
        \arrayLine{\textsc{p}, \textsc{np}, \textsc{np-complet}\\  $\leq_p$, $\equiv_p$}
        \arrayLine{\structure{\textsc{NPUniversal}} \\ est \textsc{np-complet}}
        \arrayLine{\structure{\textsc{sat}} \\ est \textsc{np-complet}}
        \arrayLine{\onlyb<1>{{\Large\structure{?}}}\only<2>{Absence de preuves \structure{courtes}\\ (Cook-Reckhow)}}
      }
    \end{tikzpicture}
  }

  \onBlock<2->[bottom=-2mm]{Théorème -- Cook-Reckhow (sous la conjecture $\textsc{np} \neq \textsc{co-np}$)}{
    Il n'existe pas de système de preuve $\langle \Phi, \Pi, \valid \rangle$ à la fois : 
    \begin{description}[Propositionnel :]
    \item[Polynomial :] les preuves sont vérifiables en temps polynomial
    \item[Propositionnel :] les propositions de $\Phi$ sont les formules en FNC
    \item[Exact :] les propositions démontrables sont les tautologies
    \item[$p$-borné :] la taille des preuves est polynomiale en celle des formules
    \end{description}
  }
  
\end{frame}

\endgroup
